\chapter{Closed cycles, splitting and external change}\label{ch:cycles}

\section{Can damage and repair mint value?}
Imagine moving water away from a useful distribution and then returning it.
The return movement may have positive EBU. If an actor reports only that
movement, it looks like a renewable source of credit. The first movement,
however, raised the represented potential and had the opposite sign. A
complete account follows both events from their actual baselines.

For an undriven sequence that begins and ends at the same complete state
$x_0=x_m$ under one fixed potential, Chapter~\ref{ch:receipts} gives
\begin{equation}\label{eq:cycle}
 \sum_{r=1}^m E_r=V(x_0)-V(x_m)=0.
\end{equation}
The theorem does not need a Gaussian. It needs the complete chain and the
same state function at every comparison. If a transition is a simultaneous
group, its exact group total is one line in the chain. The equation still
does not dictate how that line is divided among people.

At reference $(10,10)$ with scales $(2,2)$, transfer two units one way.
The new state $(8,12)$ has potential one, so the action is worth $-1$.
Transfer the same amount back and the return action is worth $+1$. The
positive restoration is real within the account; the claim of positive
\emph{net} creation from the complete loop is false.

\section{The represented state must really return}
A battery can recover its displayed charge while losing health. A tank can
recover its volume while losing quality. A vehicle can return to a depot
after using fuel. If those changes are part of the physical claim, a loop
in one displayed coordinate is not a loop in the complete state. A carrier
or boundary record must show what happened, and a potential over a fuller
state may give a different result.

Equal potential is weaker still than equal state. States $(18,2)$ and
$(2,18)$ have the same Gaussian potential but very different source stocks
and future options. Their direct finite EBU is zero; the physical state did
not return. An audit of future feasibility needs the actual coordinates,
not just their scalar value.

\section{Dividing an action cannot divide away its consequence}
From $(18,2)$, an unsplit four-unit transfer ends at $(14,6)$ and reduces
potential by twelve. Split it into two genuine sequential two-unit actions.
The first moves potential $16\to9$ and records seven. The second begins from
that new state, moves $9\to4$ and records five. Thus $7+5=12$. Calling both
pieces seven would reuse the original state twice.

This is refinement invariance of a state-function difference:
\[
 [V(x)-V(y)]+[V(y)-V(x^{\mathrm{b}})]=V(x)-V(x^{\mathrm{b}}).
\]
It does not claim that physically stopping and restarting a pump always has
the same complete effects as running it continuously. Different energy use,
heat or wear would change the complete physical transition if those carriers
were represented. The identity says that a consequence cannot be made to
disappear by changing how its record is sliced.

\section{An aggregate cycle need not restore every actor's history}
Suppose a return loop contains an action attributed to one owner and its
reverse attributed to another. The total EBU is zero, but the two historical
accounts can change by opposite amounts. A three-cell example from the
physical foundation begins at $(4,4,4)$: $B\to C$ has EBU $-1$ assigned to
$B$, and the reverse $C\to B$ has EBU $+1$ assigned to $C$. The state
returns; the owner vector changes by $(0,-1,+1)$. This is an existence
example, not a claim that every possible attribution rule has such loops.

The distinction matters because exact group closure is sometimes mistaken
for a fairness theorem. It is not. Aggregate EBU is a difference of state
values, while a person's historical attribution can depend on the path of
events and on the declared ownership rule.

\section{External change can recreate work}
If a disturbance moves a system away from its reference and an actor
restores it, positive actor EBU may recur. That does not violate
Equation~\ref{eq:cycle}: the interval contains an external event. In the
notation of Chapter~\ref{ch:receipts}, the action total is
\[
 \sum_t E_t=V(x_0)-V(x_N)+\sum_t D_{{\mathrm{ext}},t}.
\]
If repeated disturbances increase potential by four and subsequent actions
reduce it by four, the positive actor entries are balanced by real external
changes. A request or desire alone is not a physical sink. To claim a
recurring demand caused a stock change, a model must represent consumption,
transformation or boundary exchange.

Nor may an actor-created disturbance be silently renamed an external one.
The equation tells us what entries are needed; provenance and causal review
tell us whether they were classified honestly. The mathematics cannot
reconstruct an omitted event from an apparently convenient total.

\section{Changing the ruler is a third kind of event}
At the same state $(14,6)$, scales $(2,2)$ give potential four, while scales
$(4,4)$ give potential one. No water moved. The three-unit difference came
from changing the field. If old and new values are joined without a
field-change entry, intermediate potentials will not telescope as though
one ruler had been used throughout.

This is not an argument against improving a model. It is a demand that
learning about the world not be disguised as an action upon the world.
Later we ask a harder question: whether EBU values earned under different
physical fields can nevertheless share one stable physical interpretation.

Finally, no cycle equation proves recovery. On a fixed finite-mass domain a
continuous Gaussian potential is bounded even if the system stays trapped
far from its reference. A motion law, actor rule and disturbances must be
specified before attraction or long-run homeostasis can be tested. The next
chapter first examines the more immediate complication: several actions
sharing one nonlinear comparison at the same time.
